%Ex4.5, JBB

Le chlorure de sodium est parfois utilisé pour constituer des mélanges réfrigérants glace-sel.

La figure suivante présente un réseau de diagrammes d'analyse thermique pour des mélanges eau-chlorure de sodium de différentes fractions massiques (en \ce{NaCl}). Chaque courbe est obtenue en représentant l'évolution au cours du temps, lors d'un refroidissement isobare, de la température d'un système \ce{H2O}-\ce{NaCl}, dont la fraction massique en \ce{NaCl} est indiquée en haut de courbe.

\begin{center}
\centering
\begin{minipage}{0.75\textwidth}
\includegraphics[width=\textwidth]{/home/jeff/Documents/IUTOrsay/1-Cours/4p13-Alliages/Cours-TD/Exos/Binaires-Lecture/Heterogene-Ex03/figures/Fig01}
\captionof{figure}{Courbes de refroidissement thermique de divers mélange, de fraction massique $w$ en \ce{NaCl}.}
\end{minipage}
\end{center}

\Q Grâce à ce réseau, représenter le diagramme de cristallisation (binaire solide-liquide) isobare $T = f (w)$ où $w$ représente la fraction massique de NaCl dans le mélange pour les mélanges de fraction massique comprise entre $0$ et $0.25$. Une justification très soigneuse est attendue.

\Q Identifier sur le diagramme précédent les domaines limités par les courbes sachant que
l'eau et le chlorure de sodium solide forment un hydrate solide de formule \ce{NaCl,2H2O}.